\section{Two-Level CBR: Level 2 --- Abstract Case Learning}
\label{sec:level2}

Level~1 CBR (Sect.~\ref{sec:level1}) operates \textit{within} a plan. Level~2 CBR
operates \textit{across} plans: it treats the plans themselves as cases --- abstract
reasoning patterns that encode what structure of reasoning reliably solves a class of
problem.

\subsection{What a Level 2 Case Is}

A Level~2 case is a \textbf{NormCode runtime definition}: an executable plan that
specifies the reasoning structure for a task type. Plans may be authored manually,
produced through the compilation pipeline (Sect.~\ref{sec:level2:distillation}), or
developed through iterative human-AI collaboration --- the defining property is that
the result is a fully executable \texttt{.ncds} + repositories artifact, not the path
by which it was produced.

\textbf{Definition.} A \textit{Level~2 case} $P$ is a pair:
\begin{equation}
  P = (\mathcal{IR},\; \mathcal{CR}_{\text{schema}})
\end{equation}
where $\mathcal{IR}$ is the \textbf{Inference Repository}: the complete inference graph
for this plan; $\mathcal{CR}_{\text{schema}}$ is the \textbf{Concept Repository
schema}: the declared tensor axes and shapes for every concept.

\subsection{The Level 2 CBR Cycle}

\textbf{Retrieve}: a practitioner reads the \texttt{.ncn} narrative and determines
whether an existing plan (or part of it) matches the new task class.
\textbf{Reuse}: the retrieved plan is instantiated with new input values and run; all
four production-ready plans are reused in this sense on every run, each generating a
fresh set of Level~1 cases. \textbf{Revise}: when the plan does not suit the new task,
the designer modifies the \texttt{.ncds} source and recompiles; the \texttt{.ncn}
re-review verifies the revised abstract case before it generates any concrete cases.
\textbf{Retain}: the revised or newly authored plan joins the case library, which grows
through use.

\subsection{The Compilation Pipeline for NormCode Production}
\label{sec:level2:distillation}

The compilation pipeline (Derivation $\rightarrow$ Formalization $\rightarrow$
Post-Formalization $\rightarrow$ Activation) is one standardized distillation path that
produces Level~2 cases from experience and intent: \textit{Derivation} extracts
conceptual structure from natural language; \textit{Formalization} assigns flow indices
and determines execution types; \textit{Post-Formalization} assigns execution resources;
\textit{Activation} serializes the abstract case into orchestrator-ready JSON.

A notable consequence of NormCode as a reasoning medium is that \textbf{the
compilation pipeline itself is a NormCode plan.} This pipeline is a Level~2 case (the NC~Compilations plan) --- but it is not a
universal compiler. What it produces depends on the task's nature and the quality of
the input specification. For plans that involve rich domain logic, external tooling, or
complex control flow, manual authoring or human-assisted iteration is the practical
path.

An initial \texttt{.ncds} sketch of the compilation plan:

\begin{lstlisting}
<- compilation result
    <= report compilation status and output paths
    <- task specification
    <- project context
    <- specification analysis
        <= analyze specification: identify root output,
           ground inputs, and transformations
        <- ....
    <- ncds draft
        <= derive hierarchical concept structure
           from specification analysis
        <- ....
    <- formal ncd
        <= formalize draft: assign flow indices,
           sequence types, and value bindings
        <- ....
    <- post-formal ncd
        <= post-formalize: assign paradigms,
           reference shapes, and provision paths
        <- ....
    <- repositories
        <= activate post-formal plan into executable
           concept and inference repositories
        <- ....
\end{lstlisting}

Each \texttt{<=} in this initial sketch is a high-level operation whose execution is
underspecified. The compilation process itself decomposes these into concrete inferences.
This progressive decomposition is itself a Level~2 revision cycle --- improving accuracy
by giving finer inferences more precisely scoped inputs.

The recursive structure has three consequences: (1)~\textbf{Validation} --- the pipeline
is self-consistent; it can compile itself; (2)~\textbf{Improvement through revision} ---
improving the compilation pipeline is abstract case revision at the meta-level;
(3)~\textbf{Recursive retention} --- each improvement joins the case library.

\textbf{Scope and limits of automated compilation.} NC~Compilations is an initial
demonstration, not a universal solution. For plans with clear, sequential structure ---
document analysis, report generation, contract review, or compilation steps supported by
explicit provisions of scripts and prompts --- the pipeline reliably produces executable
plans within fewer than ten adjustments. Plans with complex domain logic, rich
environment interaction (e.g., PPT~Generation with its dual-format output loop), or the
compilation pipeline itself, require human or code-assistant inspection during authoring:
the structure is too open-ended for automated derivation to produce a correct plan
without iterative review. The boundary between automated and human-assisted authoring
is task complexity, not a limitation of the NormCode language.

Level~1 and Level~2 are reciprocally constitutive: every Level~1 run is an instance of
exactly one Level~2 case; experience accumulated as Level~1 cases informs revision of
Level~2 cases; the compilation pipeline is itself a Level~2 case instantiated to produce
new Level~2 cases.
